\documentclass[12pt]{article}

% --- Page Setup ---
\usepackage[paper=a4paper, left=2cm, right=2 cm, top=1.5cm , bottom=3cm]{geometry}


% \renewcommand{\baselinestretch}{1.3}  % line spacing

% --- decorations ---
\usepackage{tcolorbox}
\usepackage{xcolor}
\usepackage{fontspec}
% \selectcolormodel{gray}  % Force all color to grayscale
\usepackage{fontawesome5}
\usepackage{tabularx}
% --- Persian Support (for Persian text) ---
\usepackage{xepersian}
\settextfont{XB Zar.ttf}  % Example Persian font
\newfontfamily\XBZarBd{XB Zar Bd.ttf}
\newfontfamily\titr{XB Titre.ttf}

% --- Fancy Header and Footer ---
\usepackage{fancyhdr}
\usepackage{graphicx}
\usepackage{array}

\pagestyle{fancy}

% Header: left = placeholder for logo, right = Physician's name/clinic

% \fancyhead[R]{\raisebox{-1\height}{\includegraphics[width=5cm]{"logo.png"}}}



\fancyfoot[L]{\footnotesize
\\
  \faBuilding \space
 تهران، کارگر شمالی، بالاتر از مرکز قلب تهران، بن بست سوم، پلاک 2، طبقه 4، واحد 10
}

\fancyfoot[C]{ صفحۀ \thepage}

\fancyfoot[R]{\footnotesize
  88533271-021\space \faPhone*
  \\
   www.Drarzaghi.com  \space   \faGlobe
}

\renewcommand{\headrulewidth}{0pt}
\renewcommand{\footrulewidth}{0.4pt}

% --- Document Begins ---

\begin{document}
\begin{table}[h!]
\begin{tabular}{m{5.5cm} m{4.5cm} m{5.5cm}}
\includegraphics[width=4cm]{"<<base-dir>>/latex/compact.png"}
% \Large{\titr{دکتر سید مسعود ارزاقی}} \newline
% {\large{\XBZarBd{متخصص اعصاب و روان}}}  \newline
% {\normalsize{\XBZarBd{فلوشیپ فوق تخصصی سایکوسوماتیک}}}

&

\begin{tcolorbox}[
    colback=black!50,    % Background color
    colframe=black!50,   % Border color
    coltitle=white,     % Title text color
%     Rounded corners,   % sharp corners,
     boxrule=0.1mm,      % Border thickness
    width=\linewidth,          % Width of the box
%     height=1cm,         % Height of the box
    halign=center,      % Align the text to the center
    valign=center,      % Vertically center the text
    fonttitle=\bfseries,
    coltext=white,      % Text color
    fonttitle=\sffamily
]
\normalsize
\titr برگه‌ی توافق‌نامه
\end{tcolorbox}

&

\raggedright
\includegraphics[width=3cm]{"<<base-dir>>/static/common/logo.png"}

\end{tabular}
\end{table}
\hrule


% --- Patient Demographics Table ---
\begin{table}[h!]
\centering
    \renewcommand{\arraystretch}{1.2}  % Row spacing
     \footnotesize

     \newcolumntype{B}{>{\fontsize{8}{9.5}\selectfont\XBZarBd}p{1.5cm}}  % Bold and footnotesized column
     \newcolumntype{N}{p{3.35cm}}  % Normal column
    \begin{tabular}{B N | B N | B N }
         شماره پرونده & <<patient-id>> &
        \textbf{کد ملی} & <<ssn>> &
        \textbf{بیمه} & <<insurance>> \\
        \textbf{نام} & <<first-name>> &
        \textbf{نام خانوادگی} & <<last-name>> &
        \textbf{تاریخ تولد} & <<birthday>> \\
        \textbf{جنسیت} & <<gender>> &
        \textbf{قد} & <<height>> &
        \textbf{وزن} & <<weight>> \\
        \textbf{وضعیت تاهل} & <<marital-status>> &
        \textbf{تحصیلات} & <<education-level>> &
        \textbf{شغل} & <<occupation>>  \\

      \textbf{زبان مادری} & <<language-option>>  &
      \textbf{قومیت} & <<ethnicity>>                    &
      \textbf{گروه خونی} & <<blood-group>>  \\

      \textbf{شماره تماس} & <<phone>>             &
      \textbf{شماره واتساپ} & <<messenger1>>    &
      \textbf{شماره ایتا} & <<messenger2>>  \\

      \textbf{دست غالب} & <<dominanthand>>             &
      \textbf{آدرس} & \multicolumn{3}{>{\raggedleft\arraybackslash}p{9cm}}{<<address>>}

    \end{tabular}

\end{table}

\hrule
\vspace{1cm}
\large
\titr
توافق نامه بیمار و کلینیک

\vspace{1cm}
\normalfont \normalsize
این قرارداد فی‌مابین دکتر سید مسعود ارزاقی (که از این پس «پزشک» نامیده می‌شود) از یک سو و <<patient>> (واجد اطلاعات درج شده بالای این برگه، که از این پس «مُراجع» نامیده می‌شود) از سوی دیگر، براساس قوانین جاری کشور و اصول اخلاقی و حرفه‌ای، منعقد و لازم‌الاجرا می‌گردد. مفاد این قرارداد ناظر بر تعهدات و الزامات طرفین در ارائه و دریافت خدمات تشخیصی، درمانی و پژوهشی است.

\vspace{0.5cm}
\large \XBZarBd
ماده ۱ موضوع قرارداد
\normalsize \normalfont
ارائه خدمات پزشکی، شامل معاینه، تشخیص و انجام درمان لازم از سوی پزشک، و همکاری کامل مراجع در امور درمانی، با رعایت موارد حرفه‌ای و قانونی.

\vspace{0.5cm}
\large \XBZarBd

ماده ۲ حقوق و تعهدات مراجع
\normalsize \normalfont

\begin{enumerate}
\item
مدت زمان ویزیت پزشک برای بیماران جدید حداکثر 30 دقیقه و برای بیماران دارای پرونده حداکثر 15 دقیقه می‌باشد. در صورت طولانی‌تر شدن مدت ویزیت، زمان اضافه مشمول هزینه خواهد بود.
\item
مراجع موظف است پس از ویزیت حضوری، زمان مراجعه بعدی خود را در هماهنگی با پذیرش مشخص نماید.
\item
در صورت عدم امکان مراجعه در زمان تعیین شده، مراجع مکلف است حداقل ۴۸ ساعت پیش از وقت مقرر، موضوع را به پذیرش اطلاع داده و وقت بعدی را هماهنگ نماید. (عدم اطلاع‌رسانی ممکن است منجر به محدودیت در تعیین وقت بعدی گردد.)
\item
مراجع آگاه است که زمان ویزیت پیش از ساعت ۶ عصر روز قبل، از طریق اطلاعات تماس موجود در پرونده قطعی خواهد شد و در صورت عدم پاسخگویی، وقت به مراجع دیگری تخصیص می‌یابد.
\item
مراجع باید جهت حفظ بهداشت، نظافت مطب را رعایت کرده و با پرسنل در حفظ آرامش و نظم مشارکت نماید.
\item
در صورت وجود هر گونه نارضایتی از خدمات ارائه‌شده (پذیرش، تعیین وقت، خدمات درمانی و ...)، مراجع می‌تواند برای شماره‌ی ۰۹۰۵۸۹۶۱۸۵۰ پیام بگذارد.
\item
به منظور حفظ آرامش سایر مراجعین، در صورتی که خردسال شما بیمار یا مراجع نیست، حدالامکان از همراهی کودکان خودداری نمایید.
\item
در صورت مراجعه خارج از وقت تعیین‌شده یا ارجاع از سوی همکاران، مراجع می‌پذیرد که ممکن است تا زمان مساعد برای ویزیت صبر نماید و اولویت همواره با مراجعانی است که پیشتر وقت گرفته‌اند یا بیماران اورژانسی.
\item
مراجع موظف است زودتر از ۱۵ دقیقه به وقت تعیین‌شده وارد مطب نشود تا ازدحام در اتاق انتظار پیش نیاید.
\item
مطابق دستورالعمل‌های سازمان نظام پزشکی و پزشکی قانونی، مراجع تأیید می‌نماید که آگاهانه پزشک خود را در انتخاب روش‌های معالجه و درمان، امین دانسته و از پزشک بابت هرگونه اقدامات ضروری درمانی، برائت می‌دهد.
\item
مراجع آگاه است که هزینه‌های تشخیصی و درمانی که پس از ویزیت حضوری پرداخت شده است، در هیچ زمانی قابل استرداد نخواهد بود.
\item
مراجع تحت هیچ شرایطی نباید پشت میز پذیرش یا در اتاق‌های مطب بدون هماهنگی حضور یابد.
\item
مراجع تأیید می‌نماید که اطلاعات مربوط به پرونده و زمان و تعداد مراجعات، جزو اطلاعات شخصی و محرمانه محسوب می‌شود و هیچ یک از پرسنل، درمانگران و پژوهشگران مجاز به دراختیار قرار دادن آن‌ها به دیگران نیستند.
\newline
\XBZarBd تبصره:
\normalfont
شرایط استفاده از اطلاعات مراجع برای پژوهش های علمی در ماده ۳ تشریح شده است.
\item
مراجع با آگاهی از ماده ۶۱۸ قانون مجازات اسلامی، تعهد می‌نماید از هرگونه توهین و رفتار نامناسب با پرسنل درمانی به هنگام انجام وظیفه خودداری کرده و موجب اخلال در نظم و آرامش عمومی نشود. مصادیق توهین شامل استعمال الفاظ رکیک، تحقیر، فحاشی و اقدامات عملی (ضرب و ...) تلقی می‌شود.
\item
شماره تلفن 88533271-021 و شماره پیامرسان کلینیک ۰۹۰۵۸۹۶۱۸۵۰ برای ارسال آزمایشات و پاسخ‌دهی به سؤالات دارویی تعیین گردیده است.
\end{enumerate}


\vspace{0.5cm}
\large \XBZarBd
ماده ۳ رضایت نامه استفاده از داده های مراجع در تحقیقات
\normalsize \normalfont

\begin{enumerate}
\item
مراجع بدین‌وسیله رضایت می‌دهد که داده‌های مربوط به روند تشخیصی و درمانی وی، در چارچوب پروژه‌های پژوهشی و علمی مورد استفاده قرار گیرد، مشروط بر اینکه حریم خصوصی و اصول محرمانگی رعایت شود.
\item
تمامی اطلاعات هویتی مراجع پیش از هرگونه استفاده‌ی علمی، به‌صورت ناشناس یا کدگذاری‌شده نگهداری و پردازش خواهد شد؛ و هیچ اطلاعاتی که منجر به شناسایی وی گردد، در اختیار پژوهشگران ثالث قرار نخواهد گرفت.
\item
مراجع حق دارد در هر زمان که بخواهد، از ادامه‌ی اشتراک‌گذاری داده‌های خود در این پژوهش انصراف دهد؛ بدیهی است داده‌های پردازش‌شده یا منتشرشده تا پیش از اعلام انصراف، همچنان مورد استناد پژوهش خواهد بود.
\item
مراجع آگاه است داده‌های فوق ممکن است برای ارتقای روش‌های تشخیصی و درمانی، آموزش تخصصی، یا تحقیقات دانشگاهی به کار گرفته شود و مشارکت وی ممکن است به پیشرفت خدمات درمانی و بهره‌مندی سایر بیماران در آینده کمک نماید.
\item
به درخواست مراجعین و پس از اتمام پژوهش، خلاصه‌ای از نتایج مرتبط با پژوهش (تا حدی که منجر به افشای هویت افراد نگردد) در اختیار ایشان قرار خواهد گرفت.
\item
پزشک یا نماینده وی مسئول است اطلاعات لازم درباره‌ی پژوهش را در اختیار مراجع قرار دهد و در صورت بروز هرگونه پرسش، نگرانی یا ابهام، راه‌های ارتباطی با مسئول پژوهش را معرفی کند.
\end{enumerate}


\vspace{0.5cm}
\large \XBZarBd
ماده ۴ حقوق و تعهدات پزشک
\normalsize \normalfont

\begin{enumerate}
\item
پزشک متعهد است در ارائه خدمات تشخیصی و درمانی، اصول علمی و موازین اخلاق پزشکی را رعایت نموده و نهایت تلاش خود را در معالجه مراجع به کار گیرد.
\item
پزشک مسئول است پرونده‌ی مراجع را محرمانه دانسته و تنها در موارد مقتضی و با رعایت قوانین مربوطه، اطلاعات آن را در اختیار مراجع ذیصلاح قانونی قرار دهد.
\item
برنامه‌ی زمانی مطب مطابق با قوانین و شرایط اعلام‌شده عمل می‌کند:

\begin{itemize}
\item
واحد پذیرش: از ساعت ۱۴:۰۰ تا ۲۱:۰۰
\item
 واحد rTMS
 \footnote{تحریکِ مغناطیسیِ مکررِ فراجمجمه‌ای یا Repetitive Transcranial Magnetic Stimulation}
 و واحد LORETA
 \footnote{توپوگرافی فعالیت الکتریکی مغز با رزولوشن اندک یا Low Resolution Brain Electromagnetic Tomography}
از ساعت  ۱۰:۰۰ تا ۲۰:۰۰
\end{itemize}

\item
هرگونه تغییر در ساعت کاری، مراتب از طریق پیامک، تلفن یا اطلاعیه در مطب اعلام خواهد شد.
\end{enumerate}

\vspace{0.5cm}
\large \XBZarBd
ماده ۵ سایر شرایط
\normalsize \normalfont

\begin{enumerate}
 \item


این قرارداد صرفاً یک چارچوب کلی است و در صورت نیاز، جزئیات تکمیلی درمانی در نسخه، دستور پزشک یا فرم‌های رضایت‌نامه جداگانه قید خواهد شد.
 \item
هرگونه اختلاف در تفسیر یا اجرای مفاد این قرارداد، ابتدا از طریق مذاکره میان طرفین حل و فصل می‌گردد و در صورت عدم توافق، تابع قوانین جاری کشور خواهد بود.
\end{enumerate}

\vspace{0.5cm}
\large \XBZarBd
ماده ۶ تاریخ و امضاء
\normalsize \normalfont

این قرارداد <<today>> در محل کلینیک دکتر سید مسعود ارزاقی تنظیم و به امضای طرفین رسیده و لازمالاجراست.

\vspace{0.5cm}
% \begin{table}[h]
% %     \renewcommand{\arraystretch}{1.5} % Adjust row height
% %     \begin{tabular}{rl}
    \begin{tabularx}{\textwidth}{X X}
\XBZarBd
نام و نام خانوادگی مراجع یا قیّم قانونی وی: \vspace{2cm}
  &
  \XBZarBd
 نام پزشک:  \newline  \normalfont دکتر سید مسعود ارزاقی
 \\

\XBZarBd
 امضاء و اثر انگشت مراجع یا قیّم قانونی وی:
 &
 \XBZarBd
 امضاء و مهر پزشک یا نماینده قانونی:
   \\



    \end{tabularx}
% \end{table}





\end{document}
